\subsection{Estimates for the independent source maps}
\label{rw:app:source-estimates}

Fix every derivative order before the coefficient box and scale bound.
Use the selected jets of Section~\ref{rw:source:construction}, retaining
two extra coefficient orders for divided coefficients. We construct the
core, tail inverse and upper return on independent common domains before
imposing physical scale relations; constants may depend on the finite
order.

\begin{proof}[Core estimate in Proposition~\ref{rw:source-maps}]
Put $H=\vartheta^{-2}$ at a real base point. Complexify the independent
parameters in the product
\[
 |\widetilde\vartheta/\vartheta-1|<1/32,\quad
 |\widetilde\beta_i|<2b_0,\quad |\widetilde{\widehat\ell}|<4,
 \quad |\widetilde s|<\tfrac32e^{-H}.
\]
The integration variable $X$ remains real. With tildes suppressed in
the functions, the elementary estimates
\[
 |F-1|<5b_0,\quad
 |(\vartheta/\widetilde\vartheta)^2/F-1|<1/10,
 \quad |\widetilde\beta_c/F|<3b_0
\]
hold after fixing $b_0$ small. Integration of the real part of the
exponent of $J$ gives
\begin{equation}\label{rw:app:source-Gaussian}
 \frac9{20}HX^2-3b_0X\le\log|J(X)|
            \le\frac{11}{20}HX^2+3b_0X.
\end{equation}
In particular $|J|\le e^{3H/5}$, whereas
$|J^{-1}|\le e^{3b_0}e^{-9HX^2/20}$. Consequently
\[
 |I(X)|\le\frac{32e^{3b_0}}{31(1-5b_0)}
                    \int_0^\infty e^{-9t^2/20}\,dt<2.
\]
Rescaling the reciprocal Gaussian by $\widetilde\vartheta$ therefore
bounds $I$ uniformly as the scale vanishes.

The normalized equation can be continued through signed zeros of $U$
because its remainder contains no division by $U$. Write
$m_* =\widetilde\vartheta^2\widetilde\beta_m$,
$c_* =\widetilde\vartheta\widetilde\beta_c$,
$A_X=X/\widetilde\vartheta-c_*$ and $H_z=1+2m_*-m_*z$.
Direct substitution into \eqref{rw:source:core-equation} gives
\[
 U_X=\frac{\widetilde{\widehat\ell}}{\widetilde\vartheta FJ}
       +\mathcal R_c,\qquad U(0)=1,
\]
where
\[
 \mathcal R_c=\frac{\widetilde s}{\widetilde\vartheta}
 \left\{\frac{\widetilde{\widehat\ell} UH_z}{DF}
 -\frac{2\widetilde{\widehat\ell} U}{D}
 +JU^2\left(\frac{A_XH_z}{DF}
                    +\frac{c_*+\widetilde s\widetilde{\widehat\ell}}{D}\right)
 \right\}.
\]
This equation is holomorphic also at $\widetilde s=0$. On a
provisional tube $|U|\le16$ one has
$|z|\le24e^{-2H/5}<1/8$, $|F|>3/4$ and $|D|>1/2$.
It follows, with an absolute constant $C$, that
\[
 |\mathcal R_c|\le C|\widetilde s|H(1+|J|),\qquad
 \sup_{0\le X\le1}\left|\int_0^X\mathcal R_c\,dt\right|
                              \le CH e^{-2H/5}.
\]
For large $H$ this last quantity is less than $1/48$.
The integral equation now improves the tube to
$|U|\le1+4\cdot2+1/48<10$. Continuation on the compact
integration interval proves existence on the entire product and the
same bound for $R=U-1-\widetilde{\widehat\ell} I$.

At a real point, Cauchy radii $\vartheta/64,s/4,b_0/2,1/2$
lie in this product. They bound every weighted ordinary derivative
of $R$ by $C_NH e^{-2H/5}$, and hence every selected Euler word
by $C_Ne^{-H/3}$. The same argument applied to $I$ bounds its
selected jets. The independent neighborhoods around individual real
points also cover the boundary $s=e^{-H}$. On the real slice
$D\ge7/8-3b_0>27/32$.
On $0\le X\le\vartheta$ one has $J<2$ and $F<9/8$, so
$I(1)>4/9$. Thus $U_{\widehat\ell}(1)>1/4$, after another increase of
$H$. At $\widehat\ell=-3$ the endpoint is negative, whereas at zero it
is positive. Strict monotonicity gives the unique simple threshold;
in fact
\[
 \widehat\ell_c=-I(1)^{-1}+O(e^{-H/3}),\qquad
 U(1)=(\widehat\ell-\widehat\ell_c)G_c,\quad 1/4<G_c<3.
\]
The threshold and the factor inherit the finite parameter jets by
the implicit function theorem and its mean-value formula.

Finally, at any real zero of $z$ its derivative has the sign of
$\widehat\ell$. If $\widehat\ell<0$, all such zeros cross downwards, so a
second zero is impossible; its endpoint sign determines passage.
If $\widehat\ell>0$, a positive solution cannot cross downwards; if
$\widehat\ell=0$, uniqueness prevents a zero. These observations give
the whole-path threshold assertion, including its equality case.
\end{proof}

\begin{proof}[Independent tail and inverse estimates]
Keep the preceding complex neighborhoods and take, around each real
$v\in[7/16,9/16]$, a complex disc of radius $1/32$.
Decrease the scale radius if necessary so that
\[
 H/3<\operatorname{Re}(\widetilde v/\widetilde\vartheta^2)<2H/3,
 \qquad |\widetilde v|<3/5.
\]
Hence the exponential in \eqref{rw:source:tail-equation} satisfies
\[
 |z|\le e^{-(1-r)H/3},\quad
 |z^{-1}|\le e^{2(1-r)H/3},\quad
 \int_0^1|z|\,dr\le3/H.
\]
Writing its denominator as $X+\eta_d$, one obtains
\[
 |\eta_d|\le Cb_0\vartheta^2+C\vartheta e^{-H/3}.
\]
The second term uses $|s/z|\le(3/2)e^{-H/3}$ and therefore
holds before any physical entry is inserted.

Use $E=X^2/2$. On a provisional tube
$\operatorname{Re}X>3/4$, $|X|<2$, its equation has denominator
bounded away from zero. Integration of the equation and the preceding
integral of $|z|$ give
\[
 |E(r)-(1/2+r\widetilde v)|\le C(b_0+\vartheta^2),
 \qquad 0\le r\le1.
\]
Indeed the contributions of $\beta_BX^2+\beta_aX$ are $O(b_0)$;
those of $z,\vartheta^2\beta_m$ and $\eta_d/X$ have integrals
$O(\vartheta^2)+O(b_0\vartheta^2)+O(\vartheta e^{-H/3})$.
For small fixed $b_0$, then small $\vartheta_0$, the right side
improves the tube, for example to $\operatorname{Re}X>9/10$,
$|X|<3/2$. The branch of the square root containing $X(0)=1$
therefore continues to $r=1$.

Cauchy estimates on smaller fixed entry discs bound all entry jets,
and those on the parameter product bound the selected jets. In
particular $|T_v-1|<1/64$ after the same fixed choices. The endpoint
values at $v=7/16,9/16$ enclose $[31/32,33/32]$ strictly, since
their limits are $15/16,17/16$. The implicit function theorem on
this fixed interval, with $T_v>63/64$, supplies the inverse and its
finite mixed jets. On the real slice the original time satisfies
$r'=\vartheta(X+\eta_d)/v>0$; this is positive on the closed
$r$-interval for every fixed field and scale, so the passage takes
finite positive time.

At the same state, the difference of the tail vector fields for
$\widehat\ell$ and zero is $\widehat\ell$ times
$O(\vartheta e^{-H/3})$. Their state derivatives are bounded on
the tube. Variation of constants and analytic division in $\widehat\ell$
therefore give the divided estimate
\eqref{rw:source:tail-forcing}, including its selected jets by the
same complex neighborhoods. To divide the inverses, use
\[
 T_{\widehat\ell}^{-1}(e)-T_0^{-1}(e)
 =-\frac{T_{\widehat\ell}(T_0^{-1}(e))-e}
 {\displaystyle\int_0^1T_{\widehat\ell}'
       (T_0^{-1}(e)+t[T_{\widehat\ell}^{-1}(e)-T_0^{-1}(e)])\,dt}.
\]
The real denominator is uniformly positive; differentiation of this
identity gives the same finite jet bounds for the divided inverse.
\end{proof}

\begin{proof}[The upper passage and its endpoint layers]
Set $g(Y)=\beta_mY-1$, $f(X)=\beta_BX^2+\beta_aX$,
$G(X,Y)=X+\beta_cY$ and
\[
 Q=\frac{\lambda}{Y+\lambda},\qquad
 r_u=\frac{Y^2}{Y+\lambda}.
\]
On either endpoint layer the exact graph equation is
\begin{equation}\label{rw:app:source-upper-graph}
 \frac{dX}{dY}=\frac{(1-Q)g+Qf}{G+\alpha r_u}.
\end{equation}
Fix the layer height $d=1/8$. On the right tube $1<X<2$,
and the left tube $-2<X<-1$, $0\le Y\le d$, the denominator
has absolute value at least $7/8$. For $E=X^2/2+Y$ one has
\[
 \frac{dE}{dY}
 =\frac{Y(\beta_c+\beta_mX)+QX(f-g)+\alpha r_u}
                {G+\alpha r_u}.
\]
Since
$\int_0^d Q\,dY=\lambda\log(1+d/\lambda)$,
the energy variation on each layer is made less than $1/512$
by the fixed choices. The resulting energy bounds improve the assumed
$X$ tubes: at the most restrictive endpoint,
$X^2/2\ge3/4-d-1/512>1/2$.

In the middle $Y\ge d$, divide the field by $Y$ to obtain
\[
 \dot X=g+\lambda f/Y<-7/8,\qquad
 \dot Y=(1+\lambda/Y)G+\alpha Y.
\]
On $|X|\le2$, $d\le Y\le3/2$, the energy derivative is bounded
by $C(|\beta|+\lambda+|\alpha|)$ and the travel time by a fixed
constant. The total energy variation is less than $1/128$ after
shrinking the parameters. As long as $|X|\le1$, the energy lower
bound implies $Y>d$; on $X<-1$ the vertical derivative is strictly
negative. The arc therefore reaches the left layer transversely,
before leaving this tube. On that layer $Y$ decreases to zero and
the original vertical velocity there is $\lambda X<0$. The analogous
right-layer velocity starts positive. These signs prove that the
constructed passage is the first upper passage and has positive time.

Use \eqref{rw:app:source-upper-graph} on the two fixed $Y$-intervals.
Every $E_\lambda$ derivative of
$Q$ is $Q$ times a polynomial in $Q$; also
$E_\lambda r_u=-Qr_u$, with the same property after iteration.
Thus all selected graph-field derivatives are bounded on the tube.
For a mixed jet of order $k$, the variational equation has the form
\begin{equation}\label{rw:app:bounded-variation}
 Z_k'=AZ_k+P_k,\qquad \|A\|\le C,
\end{equation}
where $P_k$ contains bounded lower jets. Variation of constants on the
fixed layer intervals, followed by the transverse middle exit, proves
the mixed bounds by induction. Subtracting two such equations bounds
their jet differences by the integrated field-error carrier and lower
jet differences. This also proves a bounded analytic divided difference
$(\mathcal U_\alpha-\mathcal U_0)/\alpha$.
\end{proof}

\subsection{Principal maps and their differentiated errors}
\label{rw:app:source-principal}

We compare the three unforced maps with their limits in the reserved
mixed jets. This gives the divided source coefficients and their
$O(r_s)$ errors on the full contraction cube.
At $\widehat\ell=0$ the core expansion yields
\[
 v_{\pm,0}=1-\int_0^1
       \frac{t}{1+\beta_{a,\pm}t+\beta_Bt^2}\,dt+O(\vartheta^2)
\]
in every prescribed selected jet after the physical $s(\vartheta)$
is substituted. The exponential remainder absorbs the polynomial
factors introduced by $E_\vartheta s=2hs/(1-s)$.
For the unforced tail, its limiting energy equation is
\[
 E_r=v(1+2\beta_BE+\beta_a\sqrt{2E}),\qquad E(0)=1/2.
\]
The integral of the difference from the exact equation is
$O(\vartheta^2)$ on the common tube. This follows term by term
from $\int|z|\,dr=O(\vartheta^2)$ and the explicit
$\vartheta^2\beta_m,\vartheta^2\beta_c$ terms. Selected
Euler derivatives of $z$ have the form $P((1-r)/\vartheta^2)z$,
whose integrals have the same bound. Use
\eqref{rw:app:bounded-variation} on the fixed $r$-interval with this
$O(\vartheta^2)$ carrier and bounded coefficient and entry jets.
Induction gives the comparison in every reserved mixed derivative.
In the limit, $dS_\beta(X)/dr=v$, so inserting the
principal core entry gives $S_\beta(\xi_\pm)=1$.

For the upper unforced passage the limiting field is
$\dot X=-1+\beta_mY$, $\dot Y=X+\beta_cY$.
On its endpoint layers the difference of the exact graph equation
from the limiting equation is exactly $Q(f-g)/G$.
Every selected Euler derivative of this difference is bounded by
$C_NQ$ after coefficient and state derivatives. Thus the integral
carrier is $O(\lambda(1+|\log\lambda|))$.
In the middle the field error and its jets are $O(\lambda)$.
Equation~\eqref{rw:app:bounded-variation}, with these carriers on the
two fixed layers and the compact transverse middle arc, gives
\[
 |\mathcal U_0-\mathcal U_{\rm pr}|_N
             \le C_N\vartheta^2(1+|\log\vartheta|).
\]
The implicit middle exit is differentiated only at its uniformly
transverse section. Combining the three comparisons with the common
tail inverse gives the claimed $O(r_s)$ error for the unforced
shooting function on the entire contraction cube. The specified
Hadamard integrals consume two coefficient derivatives and preserve
that error. Hence they prove the error assertion for $A,C$ without
changing their normalization.

At $\beta=0$ the principal tail is $1/2+v$ and the principal upper
action is the identity; consequently $\mathcal M_{\rm pr}=\mathrm{id}$.
On a small coefficient box analyticity and the differentiated error
therefore give, at an independent entry,
\[
 |\partial_v\mathcal M-1|+|\partial_v^2\mathcal M|
 +|\partial_v^3\mathcal M|
 +h|D_{L,\mathrm{ind}}\partial_v\mathcal M|
 +h|D_{L,\mathrm{ind}}\partial_v^2\mathcal M|
 +h^2|D_{L,\mathrm{ind}}^2\partial_v\mathcal M|
 \le C(|\beta|+r_s).
\]
The unforced entry has the corresponding bounds
$h|D_Lv_{+,0}|+h^2|D_L^2v_{+,0}|\le C(|\beta|+r_s)$.
Here $D_{L,\mathrm{ind}}$ is the fixed-original-coefficient operator
\eqref{rw:source:physical-D}, with the entry independent until its
composition with the logarithmic entries below.

\subsection{Moving logarithmic entries and relative second derivatives}
\label{rw:app:source-moving-entry}

Write $D=D_L$, $\zeta=\zeta_+$ and
$w_t=v_{+,0}-t\zeta/h$. Let $\varepsilon_1$ bound the preceding
small principal-map and unforced-entry quantities. Since $Dh=1-s$,
$D^2h=s$ and $s\le1/h$,
\[
 |D(\zeta/h)|\le|D\zeta|/h+|\zeta|/h^2,
 \qquad
 |D^2(\zeta/h)|\le|D^2\zeta|/h
                     +2|D\zeta|/h^2+3|\zeta|/h^3.
\]
Uniformly in $0\le t\le1$, it follows that
\[
 \begin{aligned}
 |Dw_t|&\le\varepsilon_1/h+|D\zeta|/h+|\zeta|/h^2=:W_1,\\
 |D^2w_t|&\le\varepsilon_1/h^2+|D^2\zeta|/h
              +2|D\zeta|/h^2+3|\zeta|/h^3=:W_2.
 \end{aligned}
\]
Differentiating $S=\int_0^1\mathcal M_v(w_t)\,dt$ gives
\[
 |S-1|\le\varepsilon_1,\quad
 |DS|\le\varepsilon_1/h+\varepsilon_1W_1,
 \quad
 |D^2S|\le\varepsilon_1/h^2+2\varepsilon_1W_1/h
                      +\varepsilon_1W_2+\varepsilon_1W_1^2.
\]
The last term includes the full
$\mathcal M_{vvv}(Dw_t)^2$ contribution. The product rule and
\eqref{rw:source:log-comparison}, together with
$|\zeta|/h\le1$, imply
\[
 |D^j\{\zeta(S-1)\}|\le C\varepsilon_1|D^j\zeta|,
 \qquad j=0,1,2.
\]
For example, at order two the terms containing $(D\zeta)^2$
are bounded by $C|D^2\zeta|$; factors $|\zeta|/h$ are at
most one; the remaining terms reduce to the inequalities
$|\zeta|\le C|D\zeta|\le C^2|D^2\zeta|$.
Thus the relative estimate remains uniform as either normalized core
approaches zero within its retained domain.

Finally, the independent forced error $\mathcal E$ has all the
reserved jets bounded by $C_Ne^{-h/3}$. Two physical derivatives
after insertion of $v_\pm$ introduce at most two inverse core
factors, since each differentiation of a logarithm adds one such
factor. On the retained positive-core domain their product is at
most $s^{-4\delta_0}$. All other losses are fixed powers of $h$
from physical substitution and the leading factor $h\widehat\ell$.
Because $s=e^{-h-1+s}$ and $1/3-4\delta_0>1/4$, these losses
are absorbed in $C\widehat\ell e^{-h/4}$. Analytic division by
$\widehat\ell$ before differentiation preserves the explicit factor
$\widehat\ell$ even at zero. This proves \eqref{rw:source:forced-error}
and completes the estimates needed for the source curvature.
